471
8.4 Real-valued Functions of Several Variables
where $z_0 = f (x_0, y_0)$. Comparing equalities (8.73) and (8.74), we see that the graph
of the function is well approximated by the plane (8.74) in a neighborhood of the
point $(x_0, y_0, z_0)$. More precisely, the point $(x, y, f (x, y))$ of the graph of the func-
tion differs from the point $(x, y, z(x, y))$ of the plane (8.74) by an amount that is
infinitesimal in comparison with the magnitude $\sqrt{(x - x_0)^2 + (y - y_0)^2}$ of the dis-
placement of its curvilinear coordinates $(x, y)$ from the coordinates $(x_0, y_0)$ of the
point $(x_0, y_0, z_0)$.
By the uniqueness of the differential of a function, the plane (8.74) possessing
this property is unique and has the form
$$z = f (x_0, y_0) + \frac{\partial f}{\partial x}(x_0, y_0)(x - x_0) + \frac{\partial f}{\partial y}(x_0, y_0)(y - y_0). \quad (8.75)$$
This plane is called the tangent plane to the graph of the function $z = f (x, y)$ at the
point $(x_0, y_0, f (x_0, y_0))$.
Thus, the differentiability of a function $z = f(x, y)$ at the point $(x_0, y_0)$
and the existence of a tangent plane to the graph of this function at the point
$(x_0, y_0, f (x_0, y_0))$ are equivalent conditions.
c. The Normal Vector
Writing Eq. (8.75) for the tangent plane in the canonical form
$$\frac{\partial f}{\partial x}(x_0, y_0)(x - x_0) + \frac{\partial f}{\partial y}(x_0, y_0)(y - y_0) - (z - f (x_0, y_0)) = 0,$$
we conclude that the vector
$$\left( \frac{\partial f}{\partial x}(x_0, y_0), \frac{\partial f}{\partial y}(x_0, y_0), -1 \right) \quad (8.76)$$
is the normal vector to the tangent plane. Its direction is considered to be the direc-
tion normal or orthogonal to the surface $S$ (the graph of the function) at the point
$(x_0, y_0, f (x_0, y_0))$.
In particular, if $(x_0, y_0)$ is a critical point of the function $f (x, y)$, then the normal
vector to the graph at the point $(x_0, y_0, f (x_0, y_0))$ has the form $(0, 0, -1)$ and con-
sequently, the tangent plane to the graph of the function at such a point is horizontal
(parallel to the $xy$-plane).
The three graphs in Fig. 8.1 illustrate what has just been said.
Figures 8.1a and c depict the location of the graph of a function with respect to
the tangent plane in a neighborhood of a local extremum (minimum and maximum
respectively), while Fig. 8.1b shows the graph in the neighborhood of a so-called
saddle point.